% =============================================================================
% paper_orbitvlif.tex
%
% OrbitVLIF: Brain-Inspired Decentralized Satellite Learning for
%            On-Orbit Optical Cloud Removal
%
% IEEE conference style, single-file standalone, ~6 pages target.
%
% Module naming (deliberately distinct from prior literature):
%   OrbitVLIF                                       (was: VLIFNet)
%   MFRB  = Multi-Frequency Residual Block          (was: SRB)
%   SHAM  = Spectral-Hybrid Attention Module        (was: FSTA)
%   5QS   = 5-Level Quantized Spike                 (was: MultiSpike-4)
%   CTM   = Channel-Time Modulation                 (was: MDA)
%   AGFM  = Adaptive Gated Fusion Module            (was: GatedSkipFusion)
%   DSP   = Deep Spectral Processor                 (was: L1 Block)
%
% Build:
%   pdflatex paper_orbitvlif.tex
%   bibtex   paper_orbitvlif
%   pdflatex paper_orbitvlif.tex
%   pdflatex paper_orbitvlif.tex
%
% Figures expected under ./figures/:
%   fig1.pdf                          (TikZ — system overview)
%   fig2_modules.pdf                  (TikZ — 5QS / MFRB / SHAM zoom-ins)
%   fig3_centralized_curves.pdf       (Python — A1/A2/C2 PSNR vs ep)
%   fig4_qualitative_grid.pdf         (Python — 4×4 qualitative)
%   fig5_ablation_bars.pdf            (Python — A1 vs B1/B2/B3)
%   fig6_federated_curves.pdf         (Python — F_SNN/F_ANN/F_plain)
%   fig7_centralized_4panel.pdf       (Python — Loss + PSNR, CR1+CR2)
%   fig8_energy_bars.pdf              (Python — ANN/SNN total energy)
%   fig9_per_layer_spike_rate.pdf     (Python — per-layer firing rate)
%   fig10_per_layer_energy_paired.pdf (Python — per-layer paired energy)
% =============================================================================

\documentclass[conference]{IEEEtran}
\IEEEoverridecommandlockouts

% --- Core packages ---
\usepackage[utf8]{inputenc}
\usepackage[T1]{fontenc}
\usepackage{cite}
\usepackage{amsmath,amssymb,amsfonts}
\usepackage{algorithmic}
\usepackage{algorithm}
\usepackage{graphicx}
\usepackage{textcomp}
\usepackage{xcolor}
\usepackage{booktabs}
\usepackage{multirow}
\usepackage{hyperref}
\usepackage{url}

% --- Hyperref colour conventions ---
\hypersetup{
  colorlinks   = true,
  linkcolor    = black,
  citecolor    = blue,
  urlcolor     = blue,
}

% --- Useful macros ---
\newcommand{\E}{\mathbb{E}}
\newcommand{\R}{\mathbb{R}}
\newcommand{\D}{\mathcal{D}}
\newcommand{\Dir}{\mathrm{Dir}}
\newcommand{\OrbitVLIF}{\textsc{OrbitVLIF}}
\newcommand{\MFRB}{\textsc{MFRB}}
\newcommand{\SHAM}{\textsc{SHAM}}
\newcommand{\FQS}{\textsc{5QS}}
\newcommand{\CTM}{\textsc{CTM}}
\newcommand{\AGFM}{\textsc{AGFM}}
\newcommand{\DSP}{\textsc{DSP}}

% Path for figures
\graphicspath{{figures/}}

% =====================================================================
\begin{document}

% --- Title block ---
\title{\OrbitVLIF: Brain-Inspired Decentralized Satellite Learning for
       On-Orbit Optical Cloud Removal}

\author{
  \IEEEauthorblockN{(authors anonymised for review)}\\
  \IEEEauthorblockA{(institutions anonymised for review)}
}

\maketitle

% =====================================================================
% Abstract
% =====================================================================
\begin{abstract}
Space Computing Power Networks (Space-CPN) demand on-orbit data
processing within tight power envelopes, while optical cloud
removal---an essential preprocessing primitive for remote-sensing
applications---remains dominated by computation-heavy artificial
neural networks (ANNs) ill-suited to satellite power budgets.
We present \OrbitVLIF, a brain-inspired decentralized satellite
learning framework that unifies an energy-efficient spiking U-Net
for cloud removal with a fully on-orbit federated training pipeline.
At the model level we design four building blocks: a
\emph{Multi-Frequency Residual Block} (\MFRB) that decomposes
features into parallel temporal- and spatial-frequency paths
fused via a cross-scale gate, a \emph{Spectral-Hybrid Attention
Module} (\SHAM) jointly recalibrating temporal-amplitude and
DCT-spatial dimensions, a \emph{5-Level Quantized Spike} activation
(\FQS) that quadruples per-spike representational range while
preserving spike sparsity, and an \emph{Adaptive Gated Fusion Module}
(\AGFM) for skip connections.  At the system level we deploy the
network on a $5\!\times\!10$ Walker-Star constellation under
\textsc{FedBN}\,+\,Gossip aggregation with source-level Dirichlet
non-IID partition over the joint CUHK-CR1\,+\,CR2 cloud benchmark.
On a single 4090, \OrbitVLIF{} attains $25.37$\,dB / $0.7728$
PSNR/SSIM on CUHK-CR1 ($+0.50$\,dB over a parameter-matched
ANN U-Net at $0.57\!\times$ the parameter count) and $23.20$\,dB on
CUHK-CR2.  Ablations attribute $0.60$, $0.52$, and $0.30$\,dB drops
respectively to removing \SHAM, the dual-frequency \MFRB{} path,
and \FQS{}---both single-component ablations falling
\emph{below} the plain ANN U-Net baseline, evidence that
\OrbitVLIF's modules form a co-designed whole.  Per-image
inference energy is reduced by $1.21\!\times$\,--\,$6.19\!\times$
over the ANN baseline ($20.75$\,mJ); the bracket reflects the
multi-level-spike granularity, conservatively bounded by a
$4.6$\,pJ/MAC formula and optimistically by a $0.9$\,pJ/AC formula.
To our knowledge this is the first end-to-end demonstration of
decentralized federated learning for spiking image regression on a
satellite constellation.
\end{abstract}

\begin{IEEEkeywords}
Space computing power networks, spiking neural networks, optical
cloud removal, decentralized federated learning, neuromorphic
computing.
\end{IEEEkeywords}

% =====================================================================
% Section content lands in subsequent commits — see plan in chat.
% =====================================================================

% =====================================================================
% Section I: Introduction
% Funnel structure: macro motivation → existing-method gap → field-level
% solution (Space-CPN) → technical pivot (SNN) → our contributions.
% =====================================================================
\section{Introduction}\label{sec:intro}

\IEEEPARstart{O}{ptical} Earth-observation satellites in low-Earth
orbit (LEO) capture hundreds of gigabytes of imagery per orbit,
underpinning environmental monitoring, disaster response, and
defense applications. A substantial fraction of every captured
frame is, however, occluded by clouds and must be restored before
downstream analytics. Conventional pipelines downlink raw data to
ground stations for restoration, but the satellite-to-ground link
is bottlenecked by short visible windows ($\sim$10\,min per pass for
LEO) and narrow-band radio frequency channels, incurring multi-minute
end-to-end latency that defeats latency-critical applications such
as wildfire localisation~\cite{ref:earthobs1,ref:earthobs2}.

To bypass the downlink bottleneck, recent proposals on
\emph{Space Computing Power Networks} (Space-CPN) push computation
on-board, treating each orbital plane as a mesh of compute nodes
interconnected by inter-satellite laser links
(ISLs)~\cite{ref:spacecpn1,ref:spacecpn2}. Optical cloud removal
becomes a canonical Space-CPN primitive: it is needed by every
downstream task, and its data is local to individual satellites.

\textbf{Yet the on-board energy budget is severe.} LEO platforms
draw $\sim$100--200\,W from solar arrays charged during sunlit arcs
and depleted on eclipse; dense artificial neural network (ANN)
inference and training compete with attitude control, payload, and
downlink for this shared budget. On 45\,nm CMOS, multiply-accumulate
(MAC) operations dominate ANN cost at $\sim$4.6\,pJ each, which is
prohibitive for the multi-G\,MAC convolutional or transformer
backbones used by recent cloud-removal SOTA~\cite{ref:cuhkcr,
ref:darkir}.

\textbf{Spiking neural networks (SNNs) offer a structurally
different operating point.} Activations are encoded as sparse,
event-driven binary or multi-level spikes; silent neurons skip the
multiply step entirely, and a 45\,nm accumulate-only operation
costs $\sim$0.9\,pJ~\cite{ref:horowitz}. SNNs have been shown to
match ANN accuracy within $\sim$3\,\% on classification benchmarks at
roughly $10\!\times$ lower per-layer energy~\cite{ref:flsnn}. Recent
work has applied SNNs to image deraining~\cite{ref:esdnet} and
classification in satellite federated learning~\cite{ref:flsnn},
yet \emph{pixel-level cloud removal under decentralized on-orbit
training has not been demonstrated}.

The closest precedent, the FLSNN framework of Yang \emph{et al.}~\cite{ref:flsnn},
established the feasibility of decentralized federated SNN training
on a Walker-Star constellation for the EuroSAT classification task.
Three concrete gaps prevent its direct application to operational
on-orbit cloud removal: (i) \emph{task-type gap} --- pixel
regression has different gradient structure (dense, low-frequency,
residual-dominated) from logit-driven classification; (ii)
\emph{architectural gap} --- FLSNN's small spiking-CNN /
spiking-ResNet backbones lack the U-Net structure, frequency-aware
attention, and multi-level activation that recent SNN
image-restoration work~\cite{ref:fsta-snn,ref:darkir} has shown
essential for high-fidelity reconstruction; (iii)
\emph{non-IID gap} --- satellite data heterogeneity is naturally
driven by source (sensor, atmosphere, look angle) rather than by
class label, calling for a source-level partition formulation that
is absent from FLSNN's class-Dirichlet protocol.

To close these three gaps, we present \OrbitVLIF, a co-designed
spiking image-restoration network and decentralized federated
learning system tailored for on-orbit cloud removal. Our
\textbf{contributions} are:

\begin{itemize}
\item \textbf{We propose} a quantized spiking U-Net for cloud
  removal that integrates four novel building blocks: a
  \emph{Multi-Frequency Residual Block} (\MFRB) decomposing
  features into temporal- and spatial-frequency parallel paths
  fused via a cross-scale gate, a \emph{Spectral-Hybrid Attention
  Module} (\SHAM) jointly recalibrating temporal-amplitude and
  DCT-spatial dimensions, a \emph{5-Level Quantized Spike} (\FQS)
  activation that quadruples per-spike representational range
  while preserving spike sparsity, and an \emph{Adaptive Gated
  Fusion Module} (\AGFM) for content-adaptive skip connections.
  To our knowledge, \OrbitVLIF{} is the first SNN architecture
  explicitly designed for optical cloud removal.

\item \textbf{We provide} the first end-to-end demonstration of
  decentralized federated learning for spiking image regression on
  a satellite constellation. We deploy \OrbitVLIF{} on a
  $5\!\times\!10$ Walker-Star configuration under
  \textsc{FedBN}\,+\,Gossip aggregation with \emph{source-level}
  Dirichlet non-IID partition over the joint
  CUHK-CR1\,+\,CR2~\cite{ref:cuhkcr} cloud benchmark, and compare
  three matched-protocol backbones (SNN \OrbitVLIF, ANN
  \OrbitVLIF, plain ANN U-Net).

\item \textbf{We disclose} per-image inference energy with a
  transparent dual-bound estimate that \emph{explicitly} brackets
  the cost of the multi-level \FQS{} quantization: a conservative
  upper bound at the ANN MAC formula (4.6\,pJ/MAC) and an
  optimistic lower bound at the binary AC formula (0.9\,pJ/AC).
  This honest bracket avoids the common pitfall of reporting only
  the optimistic AC formula on networks that use multi-level spikes.

\item \textbf{We conduct} extensive evaluation over $11$ training
  runs (centralized SOTA, three-way component ablation, federated
  three-backbone comparison, and Dirichlet $\alpha$-sensitivity)
  showing that \OrbitVLIF{} attains $25.37$\,dB / $0.7728$
  PSNR/SSIM on CUHK-CR1 with only $2.30$\,M parameters and a
  $1.21\!\times$\,--\,$6.19\!\times$ inference-energy reduction over
  the ANN baseline.
\end{itemize}

The remainder of the paper is organized as follows.
Section~\ref{sec:sysmodel} formalises the constellation, federated
learning problem, and source-level Dirichlet non-IID partition.
Section~\ref{sec:framework} details the \OrbitVLIF{} architecture
and decentralized training pipeline. Section~\ref{sec:experiments}
reports centralized, ablation, federated, and energy
results. Section~\ref{sec:conclusion} concludes.

% =====================================================================
% Section II: System Model lands in subsequent commits.
% =====================================================================

\end{document}
